\documentclass[12pt,a4paper]{amsart}
%\input xy
%\xyoption{all}

\DeclareMathAlphabet{\mathpzc}{OT1}{pzc}{m}{it}
\usepackage{amsfonts,amsmath,amssymb,indentfirst,mathrsfs,amscd}
\usepackage[mathscr]{eucal}
%\usepackage[active]{srcltx}




\setlength{\oddsidemargin}{0.4in}
\setlength{\evensidemargin}{0.4in}
\setlength{\textwidth}{5.5in}
\setlength{\textheight}{8.8in}
\setlength{\marginparwidth}{0.8in}
\addtolength{\headheight}{2.5pt}

\renewcommand{\thefootnote}{}

\begin{document}
\thispagestyle{empty}
\vspace*{-2cm}

 
\noindent
\begin{center}\Large\bfseries\textsf{\'Algebra 1 (grado en f\'isica)\\   curso 2012-13, primer semestre\\
	grupo 516}
\end{center}
\hrulefill{}
 
\noindent
{\sffamily    
Profesor:} Marco Zambon. \\
%{\sffamily Despacho:} 01.17.AU.103\\
{\sffamily Tutor\'ias:} por cita previa, en el despacho 407 del departamento de matem\'aticas (modulo 17). \\
{\sffamily P\'agina web del curso:} \texttt{http://www.uam.es/marco.zambon/WS12.html}

\noindent\hrulefill{}  

\noindent
{\sffamily Horario:} Lunes a Jueves, 17:30-18:20. Aula  01.11.AU.201-4.   \\ 
\noindent
{\sffamily  Fecha de examenes:} \\ 
%examen parcial: a decidir\\
ejercicios realizados en clase: 26 de octubre 2012 (viernes) y
 30 de noviembre 2012 (viernes), en clase.\\ 
examen final: 10 de enero de 2013 (jueves, tarde)\\
convocatoria extraordinaria:  17 de junio de 2013 (lunes, ma\~nana) 

\noindent\hrulefill{}  

\noindent
{\sffamily  Sistema de evaluaci\'on:}  (Ver la Guia Docente para m\'as detalles.)\\
La nota se determina a partir del siguiente promedio:\\
entrega de ejercicios realizados en clase (30 por ciento),\\
examen final (70 por ciento).

Los enunciados de los ejercicios realizados en clase consisteran de una selecci\'on de las hoja   para casa.

\noindent\hrulefill{}  

\noindent
{\sffamily  Hojas para casa:}\\
Cada  semana habr\'a una hoja de ejercicios para casa, que discutiremos durante la hora de pr\'acticas la semana siguiente. No ser\'a necesario entregar los ejercicios.
La participaci\'on durante las horas de pr\'acticas puede contribuir a subir la nota hasta un 5 por ciento.

%\noindent
%{\sffamily  Tareas para casa:} Cada 2 semanas habr\'a una hoja de ejercicios para casa, que discutiremos en clase en la semana siguiente. No ser\'a necesario entregar los ejercicios.

%\noindent\hrulefill{}  


\noindent\hrulefill{}  

 \noindent
{\sffamily Programa orientativo del curso:} \\
 \begin{itemize}
\item [I)] {\bf Introducci\'on.} Espacios vectoriales y las aplicaciones lineales entre ellos. Matrices y determinantes.\\
{\bf Los n\'umeros complejos.} Definici\'on y ejemplos. Teorema fundamental del \'algebra. Operaciones de los n\'umeros complejos y sus propiedades.\\
{\bf Espacios vectoriales.} Repaso del plano y del espacio reales. Espacios vectoriales. Bases y dimensi\'on. Productos escalares. Identidad del paralelogramo. Desigualdad de Cauchy-Schwarz. Normas y distancias. Minimos cuadrados; recta de regresi\'on.
\item [II)] {\bf Matrices y aplicaciones lineales.} Definici\'on y ejemplos de aplicaciones lineales. Relaci\'on con las matrices. Estudio de las operaciones de matrices. Composici\'on de aplicaciones lineales y producto de matrices. Aplicaciones lineales y matrices invertibles. Isomorfismos. Subespacios.
\item [III)] {\bf Sistemas de ecuaciones lineales.} M\'etodo de Gauss, y de Gauss-Jordan, para resolver sistemas de ecuaciones lineales. Estudio de las matrices que intervienen en la expresi\'on de un sistema de ecuaciones lineales y en su resoluci\'on. Pivotes. Caracterizaci\'on de las matrices invertibles como producto de matrices elementales. C\'alculo de la inversa de una matriz por el m\'etodo de Gauss-Jordan. Soluci\'on general de un sistema lineal. Subespacios fundamentales de una matriz. Rango de una matriz. Representaci\'on de una aplicaci\'on lineal en bases arbitrarias. Cambio de bases.
\item [IV)] {\bf Determinantes.} Definici\'on y propiedades de los determinantes. Existencia y unicidad. Expansi\'on por cofactores. Menores y rango. Caracterizaci\'on de las matrices invertibles por su determinante.
\item [V)] {\bf Introducci\'on a la Teor\'ia espectral.} Definici\'on autovalores y autovectores. Diagonalizaci\'on. Enunciado del Teorema Espectral para matrices reales y sim\'etricas.
\end{itemize}


\noindent\hrulefill{}  

 \noindent
{\sffamily Bibliograf\'ia:} \\
Ver la Guia Docente para una lista completa de libros aconsejados. La referencia principal ser\'a:


\begin{itemize}
\item \emph{Sergei Treil. Linear algebra done wrong.} Disponible en la p\'agina web del curso, \texttt{http://www.uam.es/marco.zambon/WS12.html}. \end{itemize}

  
Los bloques I-II corresponden al cap\'itulo 1 del libro,  y los bloques III, IV, V corresponden respectivamente a los cap\'itulos 2,3,4.

  

  
  
\end{document}
